\section{Costos, ingresos y viabilidad inicial}

\subsection{Criterio de estimación}

Los montos se expresan en dólares estadounidenses porque los proveedores cobran
en esa moneda y se convierten a colones con el tipo de cambio de referencia de
venta de CRC 456,94 por USD consultado el 1 de julio de 2026
\parencite{bccr_tipo_cambio}. Las tarifas pueden cambiar; por ello, deben
actualizarse antes de una decisión de compra.

Se distinguen tres conceptos: desembolso actual, aporte en especie y costo
comercial proyectado. Las horas de trabajo no se convierten en cero por no existir
salario; su valoración queda pendiente hasta registrar tiempo y acordar una tarifa.

\subsection{Costos actuales del prototipo}

\begin{table}[H]
\caption{Desembolso mensual actual del prototipo}
\begin{tabularx}{\textwidth}{>{\bfseries}p{3.2cm}r r X}
\toprule
Concepto & USD & CRC aprox. & Tratamiento \\
\midrule
ChatGPT Plus & 20 & 9\,139 & Gasto real de apoyo al desarrollo; la API se cobra por separado \parencite{openai_plus}. \\
Vercel Hobby & 0 & 0 & Plan actual de demostración, no presupuesto comercial \parencite{vercel_pricing}. \\
Supabase Free & 0 & 0 & Plan actual con límites y pausas por inactividad \parencite{supabase_pricing}. \\
Dominio & 0 & 0 & No adquirido. \\
\midrule
\textbf{Total efectivo} & \textbf{20} & \textbf{9\,139} & No incluye equipo, internet ni trabajo. \\
\bottomrule
\end{tabularx}
\caption*{\textit{Nota}. Conversión a CRC redondeada al colón más cercano.}
\end{table}

Vercel limita el plan Hobby al uso personal o no comercial
\parencite{vercel_terms}. En consecuencia, el costo cero del prototipo no puede
usarse como presupuesto definitivo para vender el servicio.

\begin{table}[H]
\caption{Aportes en especie pendientes de valorar}
\begin{tabularx}{\textwidth}{>{\bfseries}p{3.3cm}p{3cm}X}
\toprule
Recurso & Método de valoración & Estado \\
\midrule
Computadoras y teléfonos & Depreciación o costo de uso proporcional. & Equipo existente; monto pendiente. \\
Internet y electricidad & Proporción atribuible al proyecto. & Aporte familiar; monto pendiente. \\
Trabajo del equipo & Horas registradas por tarifa acordada. & Horas y tarifa pendientes. \\
Tutoría y apoyo educativo & Horas de acompañamiento. & Aporte institucional, no facturado. \\
\bottomrule
\end{tabularx}
\caption*{\fuentepropia}
\end{table}

\subsection{Escenario comercial conservador}

El escenario incorpora servicios adecuados para iniciar una prueba comercial sin
afirmar que sean suficientes para cualquier volumen. Vercel Pro parte de USD 20
mensuales y Supabase Pro de USD 25 \parencite{vercel_pricing,supabase_pricing}.
Para un proceso persistente se reserva una base de USD 5 en Railway, cuyo cobro
depende del consumo \parencite{railway_pricing}. Upstash dispone de un nivel
gratuito limitado para Redis \parencite{upstash_pricing}.

\begin{table}[H]
\caption{Infraestructura mensual comercial de referencia}
\label{tab:costos-comerciales}
\begin{tabularx}{\textwidth}{>{\bfseries}p{3.5cm}r r X}
\toprule
Concepto & USD & CRC aprox. & Supuesto \\
\midrule
Vercel Pro & 20 & 9\,139 & Frontend y servicios web según consumo incluido. \\
Supabase Pro & 25 & 11\,424 & Base de datos, almacenamiento y respaldos del plan. \\
Worker/cron & 5 & 2\,285 & Base de consumo; validar recursos reales. \\
Redis & 0 & 0 & Nivel gratuito inicial sujeto a límites. \\
\midrule
\textbf{Subtotal infraestructura} & \textbf{50} & \textbf{22\,847} & Antes de dominio, impuestos y excedentes. \\
Herramienta de desarrollo & 20 & 9\,139 & ChatGPT Plus, separable de la operación. \\
\midrule
\textbf{Total conservador} & \textbf{70} & \textbf{31\,986} & Base usada para punto de equilibrio de caja. \\
\bottomrule
\end{tabularx}
\caption*{\textit{Nota}. Estimación propia con tarifas públicas consultadas el 1 de julio de 2026.}
\end{table}

No se incluyen dominio, comisiones de pago, formalización, impuestos, soporte
profesional, salario ni contingencia. Esos rubros deben cotizarse antes de iniciar
ventas. El escenario tampoco garantiza disponibilidad o cumplimiento normativo.

\subsection{Escenarios de ingreso}

Los precios se diseñan para preguntar por disposición de pago y no como resultado
de una encuesta. La Tabla~\ref{tab:ingresos} combina planes para mostrar
sensibilidad financiera.

\begin{table}[H]
\caption{Escenarios mensuales ilustrativos de ingreso}
\label{tab:ingresos}
\begin{tabularx}{\textwidth}{>{\bfseries}p{2.3cm}X r r}
\toprule
Escenario & Mezcla hipotética & Ingreso CRC & Flujo sobre CRC 31\,986 \\
\midrule
Bajo & 2 gimnasios básicos. & 30\,000 & -1\,986 \\
Medio & 2 básicos y 2 intermedios. & 80\,000 & 48\,014 \\
Alto & 2 básicos, 3 intermedios y 1 avanzado. & 145\,000 & 113\,014 \\
\bottomrule
\end{tabularx}
\caption*{\textit{Nota}. Escenarios de análisis, no pronósticos de ventas. No deducen trabajo, impuestos ni otros costos pendientes.}
\end{table}

\subsection{Punto de equilibrio}

El punto de equilibrio de caja se calcula como costos fijos mensuales divididos
entre el ingreso promedio por cliente. Con CRC 31\,986 de costo conservador, se
requieren tres clientes si todos adquieren el plan básico de CRC 15\,000, o dos
clientes si el ingreso promedio alcanza CRC 25\,000. El resultado se redondea
hacia arriba porque no puede adquirirse una fracción de cliente.

Este cálculo no representa rentabilidad integral. Una vez valoradas las horas,
se aplicará:

\begin{quote}
Costo económico mensual = costos de caja + (horas mensuales $\times$ tarifa por hora).
\end{quote}

El precio definitivo deberá cubrir costo económico, riesgo y margen razonable.
Si la encuesta no respalda los rangos propuestos, se deberá reducir alcance,
ajustar infraestructura o revisar el segmento antes de vender.

\subsection{Estrategia de sostenibilidad}

La primera meta no es maximizar funciones, sino conseguir un piloto que permita
medir uso y soporte. Se recomienda limitar personalizaciones, controlar gastos,
definir respaldos y revisar mensualmente ingreso recurrente, cancelaciones,
morosidad, costo de infraestructura y horas de atención.
